\documentclass[11pt]{amsart}
\usepackage[T1]{fontenc}\usepackage{lmodern}
\usepackage[a4paper,margin=25mm]{geometry}
\usepackage{amsmath,amssymb,mathtools,microtype}
\usepackage[hidelinks]{hyperref}\usepackage{xurl}
\newtheorem{theorem}{Theorem}\newtheorem{proposition}[theorem]{Proposition}
\newcommand{\Z}{\mathbb Z}\newcommand{\Q}{\mathbb Q}
\title{Two exact positive-witness tests for Erd\H{o}s--Straus divisor boxes}
\author{The Clankers research workbench}\date{12 September 2026}
\begin{document}
\begin{abstract}
We retain all multiplicities in a signed divisor box and derive a sharp
integer collision threshold for meeting the combined exterior/middle target.
It forces a marked witness at $p=1201$ where the ordinary Cauchy--Schwarz
threshold does not. Separately, an exact largest-divisor formula sharpens the
cofactor step in an existing angular construction. The bounds are unconditional
sufficient criteria with explicit witness inverses, not hypotheses asserted
for all primes. Standard finite inequalities, the ES channels, and the
previous anchor construction are not claimed as new; historical priority of
these particular packaged refinements is unassessed.
\end{abstract}
\maketitle

Let $p$ be a prime $1\pmod4$, let $p/4<a<3p/4$, and put $R=4a-p$.
Every $u\mid a^2$ gives the exterior and middle gates
\[
 R\mid4u+1\quad(E),\qquad R\mid u+a\quad(M).
\]
With $g=\gcd(a,u)$, set $h=g^2/u$, $r=u/g$, $s=a/g$.
Then $a=hrs$, $u=hr^2$, $\gcd(r,s)=1$, and the positive ordered witnesses are
\[
 (hrs,hs\kappa,phr\kappa),\ \kappa=(pr+s)/R\quad(E),
\]
\[
 (hrs,phs\lambda,phr\lambda),\ \lambda=(r+s)/R\quad(M).
\]
Integrality follows from the indicated gate. These are the inherited
Yamamoto/Type-I--II coordinates \cite{y,et}.

\section{The integer collision threshold}
Write $a=\prod q^{e_q}$ and let $\mu(g)$ count the vectors
$\beta\in\prod_q[-e_q,e_q]\cap\Z^{\omega(a)}$ satisfying
$\prod q^{\beta_q}=g$ in $G=(\Z/R\Z)^\times$. Set
\[
 N=\prod_q(2e_q+1),\quad E=\sum_{g\in G}\mu(g)^2,
 \quad\mathcal T=\{-1,-p,-p^{-1}\}\subset G.
\]
Coincident targets are counted once. The inversion involution gives
$\mu(g)=\mu(g^{-1})$; both channels fail exactly when
$\mu$ vanishes on $\mathcal T$.

\begin{theorem}
Let $M=|G|-|\mathcal T|>0$, write $N=Mq+r$ with $0\le r<M$, and put
$L=Mq^2+r(2q+1)$. If $E<L$, one of the two gates has an actual divisor hit.
The energy can be computed without discarding any exponent budget by
\[
 E=\sum_{\delta_q=-2e_q}^{2e_q}
 \mathbf1_{\prod q^{\delta_q}=1\ (R)}
 \prod_q(2e_q+1-|\delta_q|).
\]
\end{theorem}
\begin{proof}
If there is no hit, $N$ occurrences occupy $M$ allowed bins. Whenever two
counts differ by two or more, moving one occurrence from the larger to the
smaller strictly lowers their sum of squares. Hence the minimum has $r$
counts $q+1$ and $M-r$ counts $q$, proving $E\ge L$.
The displayed energy formula counts ordered pairs of exponent vectors of
equal image, grouped by their difference. In coordinate $q$ a difference
$\delta_q$ occurs in $2e_q+1-|\delta_q|$ ordered pairs.

For the inverse, take an actual vector $\beta$ in a target fibre. If its
image is $-1$, use $M$; if it is $-p^{-1}$, use $E$; if it is only $-p$,
first negate every exponent and use $E$. In all cases
$u=\prod q^{e_q+\beta_q}$ has exponents in $[0,2e_q]$. The gate and the
ordered denominator formulas then hold exactly.
\end{proof}

There is a further exact parity improvement. If $1\in\mathcal T$, the
zero vector is a hit. Otherwise let $s$ count the allowed nonidentity
involutions and $r$ the allowed unordered inverse pairs. A no-hit histogram
has identity count $1+2x_0$, involution counts $2x_i$, and pair counts $y_j,y_j$.
Consequently its exact relaxed energy minimum is
\[
 \min_{x_0+\cdots+x_s+y_1+\cdots+y_r=(N-1)/2}
 \left((1+2x_0)^2+4\sum_{i=1}^sx_i^2+2\sum_{j=1}^ry_j^2\right).
\]
It is computed by taking the required number of least marginal costs from
$8,16,24,\ldots$ (identity), $4,12,20,\ldots$ (each involution),
and $2,6,10,\ldots$ (each inverse pair). The increasing marginal sequences
make every selected prefix feasible; an exchange proves optimality.

At $p=1201,a=306,R=23$, the actual box is
$[-1,1]\times[-2,2]\times[-1,1]$ on primes $2,3,17$.
The energy is $107$, the target set is $\{9,18,22\}$, and
$N=45$, $M=19$. Ordinary Cauchy--Schwarz does not force a hit:
$19\cdot107=2033>45^2=2025$.
But the integer threshold is
\[
 12\cdot2^2+7\cdot3^2=111>107.
\]
The inverse gives $u=17$, $E$, and
$(306,16218,1082101)$. This is a better sufficient criterion, not a claim
of new numerical coverage. The full multiplicity table is in the certificate.

\section{A sharper angular cofactor step}
Let $H\le(\Z/R\Z)^\times$ have index two, with $-1\notin H$.
Choose an actual divisor $A\mid a$ whose prime factors lie in $H$ and
whose pairs $nd\mid A$, $\gcd(n,d)=1$, represent every $v=n/d\in H$.
Define the attained rational costs
\[
 c_A(v)=\min_{n/d=v\ (R)}d/n,
 \qquad B_A=\max_{v\in H}c_A(v),\qquad b=a/A.
\]
This hypothesis concerns a finite, explicitly checked anchor; it is not
asserted for an arbitrary shell or unknown prime. It is the exact anchor
notion from the preceding workbench angular construction \cite{anchor}.

\begin{theorem}
The greatest available divisor of $b$ outside $H$ is
\[
 T_H(b)=b\quad(b\notin H),\qquad
 T_H(b)=b/q_*\quad(b\in H\text{ with an outside prime factor}),
\]
where $q_*$ is the least such prime; it is zero when no outside factor exists.
If $T_H(b)>8B_A$, the shell has a middle witness with $0<r/s<1/8$.
The exact anchor-wise criterion is
$\max_{t\mid b,\ t\notin H}t/c_A(-t^{-1})>8$.
\end{theorem}
\begin{proof}
When $b\in H$ and $t\notin H$, the complement $b/t$ is outside and must
contain an outside prime factor. Hence $b/t\ge q_*$. Equality is attained
by $t=b/q_*$. The other cases are immediate. This uses one actual prime
occurrence even when $q_*$ divides $b$ repeatedly.

For any chosen outside $t$, take a minimizing pair $(n,d)$ at $-t^{-1}$.
Then $R\mid d+tn$ and $dtn\mid a$. Put
\[
 g=\gcd(d,tn),\quad r=d/g,\quad s=tn/g,\quad
 h=ag^2/(dtn),\quad u=ad/(tn),\quad\lambda=(d+tn)/(gR).
\]
All are positive integers, $g$ is a unit modulo $R$, $a=hrs$, $u=hr^2$,
and $\gcd(r,s)=1$. The middle gate holds and $r/s=c_A(-t^{-1})/t$.
The largest-divisor sufficient condition and exact criterion follow.
For a target with the anchor retained, all inverse pairs satisfy
$t=ds/(nr)\mid b$; then $r\mid d$, $g=d/r$, and the original packet is
recovered. Thus the finite fibres are explicit.
\end{proof}
The source Koide lower boundary exceeds $1/8$, as is already certified by
$\sqrt3>1.732$, $\sqrt2<1.415$ and $317/2415>1/8$.
The coarse earlier threshold $b>64B_A^2$ is replaced, when $b\in H$,
by the factor-sensitive bound $b>8B_Aq_*$.

For $R=19,A=35$, the nine residue costs (residue, cost) are
\[
\begin{array}{c|rrrrrrrrr}
v&1&4&5&6&7&9&11&16&17\\\hline
c_A(v)&1&5&1/5&35&1/7&5/7&7&1/35&7/5
\end{array}
\]
These values follow by listing the nine coprime pairs $nd\mid35$.
At the prime $p=558721$, $a=139685=35\cdot3991$,
$b=13\cdot307\in H$, $q_*=13$, and $307>280$ but $3991<78400$.
The minimizing pair $(n,d)=(1,35)$ yields
\[
 (R,u,h,r,s,\lambda)=(19,15925,13,35,307,18),
\]
with denominators $(139685,40137399198,4575924990)$.
No claim is made that the resulting prime family was previously uncovered.

\paragraph{Scope.}
Both theorems construct actual witnesses when their proved inequalities hold.
A counterexample to ES would have to violate every sufficient test in every
full shell. No theorem that one of these tests always succeeds is assumed.
The standard-library verifier recalculates the energies, their full fibres,
the symmetric minima, all anchor pairs, and the explicit witness returns.

\begin{thebibliography}{9}
\bibitem{y} Yamamoto, K. (1965). On the Diophantine equation
$4/n=1/x+1/y+1/z$. \emph{Mem. Fac. Sci. Kyushu Univ. A, 19}(1), 37--47.
Lemma 2, equation (4), p.~38.
\bibitem{et} Elsholtz, C., \& Tao, T. (2013). Counting the number of solutions
to the Erd\H{o}s--Straus equation on unit fractions. \emph{J. Aust. Math.
Soc., 94}(1), 50--105. \S2.
\bibitem{anchor} The Clankers research workbench (2026).
\emph{Pointwise selector continuation}, section ``Weighted angular forcing''.
The supplied predecessor \texttt{es\_pointwise\_continuation\_bundle.zip}
retains the original finite anchor and its complete middle-state inverse.
The largest-cofactor step and integer collision refinement are the
present continuation. No historical-priority claim is made.
\end{thebibliography}
\end{document}
